% 本书原创：平面反射与球面反演的二维截面示意。
\begin{figure}[htbp]
\centering
\begin{tikzpicture}[x=1cm,y=1cm,>=stealth,every node/.style={font=\small},line width=.65pt]
\begin{scope}
\fill[AnalysisBlue!7] (0,0) circle(1.55);
\draw[AnalysisBlue] (0,0) circle(1.55);
\draw[->,gray] (-1.85,0)--(3.6,0);
\fill (0,0) circle(1.2pt) node[below left]{\(0\)};
\fill[AnalysisBlue] (.6975,0) circle(2pt) node[below]{\(x\)};
\fill[AnalysisBlue] (3.4444,0) circle(2pt) node[below]{\(x^*\)};
\fill (0.775,1.34234) circle(1.5pt) node[above]{\(y\)};
\draw[dashed,gray] (.6975,0)--(.775,1.34234)--(3.4444,0);
\node at (.6,2.1){球面反演};
\node at (.5,-2.0){\(R=1,\quad |x|=0.45,\quad |x^*|=1/0.45\)};
\end{scope}
\begin{scope}[xshift=8cm]
\fill[AnalysisBlue!7] (-1.7,0) rectangle(2.2,1.65);
\draw[gray] (-1.85,0)--(2.35,0) node[right]{边界};
\draw[dashed,gray] (0,-1.2)--(0,1.2);
\fill[AnalysisBlue] (0,1.2) circle(2pt) node[above]{\(x=(x',t)\)};
\fill[AnalysisBlue] (0,-1.2) circle(2pt) node[below]{\(x^*=(x',-t)\)};
\fill (1.8,0) circle(1.5pt) node[above right]{\(y\)};
\draw[dashed,gray] (0,1.2)--(1.8,0)--(0,-1.2);
\node at (.25,2.1){平面反射};
\node at (.25,-2.0){\(|x-y|=|x^*-y|\)};
\end{scope}
\end{tikzpicture}
\caption[边界配平中的两种像点]{边界配平中的两种像点。左图是单位球的二维截面，球面距离相差固定比例；右图的边界是平面，两条距离相等。像点均位于区域外，基本解的系数须分别按球面反演与平面反射公式选取。}
\label{b2:fig:13-images}
\end{figure}
